\documentclass[compress,10pt]{beamer}

%Load the myriad packages
\usepackage{amssymb,amsmath}
\usepackage{textcomp}
\usepackage{graphicx}
\usepackage{tikz}
\usetikzlibrary{spy,arrows}
\usepackage[detect-all]{siunitx}
%\usepackage[numbers, super]{natbib}
\usepackage{grffile} %spaces in file names
\usepackage{parskip}
%\usepackage[T1]{fontenc} %for sc and bf
%\usepackage{times}
\usepackage{wasysym}
\usepackage{bigstrut}
%\usepackage{hyperref}
\usepackage{float}
% Load some optional sub-parts of PGF
\usetikzlibrary{decorations.pathmorphing}
\usetikzlibrary{positioning}
\usetikzlibrary{calc}
\usetikzlibrary{shapes.geometric}
\usepackage{pgfplots}
%\usepackage[no-math]{fontspec}
%\usepackage{xltxtra}
%\usepackage{xunicode}
%\defaultfontfeatures{Mapping=tex-text}
%%\setsansfont[Mapping=tex-text]{Optima}
%\setsansfont[Mapping=tex-text]{Helvetica Neue}
% Optional for code samples
\usepackage{listings}
%
%singular
\newcommand{\fref}[1]{Fig.~\ref{fig:#1}}
\newcommand{\Fref}[1]{Figure~\ref{fig:#1}}
\newcommand{\eref}[1]{Eq.~(\ref{eq:#1})}
\newcommand{\Eref}[1]{Equation~(\ref{eq:#1})}
\newcommand{\tref}[1]{Table~\ref{tab:#1}}
%plural
\newcommand{\frefs}[2]{Figs.~\ref{fig:#1} and \ref{fig:#2}}
\newcommand{\Frefs}[2]{Figures~\ref{fig:#1} and \ref{fig:#2}}
\newcommand{\erefs}[2]{Eqs.~(\ref{eq:#1}) and (\ref{eq:#2})}
\newcommand{\Erefs}[2]{Equations~(\ref{eq:#1}) and (\ref{eq:#2})}
\newcommand{\trefs}[2]{Tables~\ref{tab:#1} and \ref{tab:#2}}
%range
\newcommand{\frefss}[2]{Figs.~\ref{fig:#1} - \ref{fig:#2}}
\newcommand{\Frefss}[2]{Figures~\ref{fig:#1} - \ref{fig:#2}}
\newcommand{\erefss}[2]{Eqs.~(\ref{eq:#1}) - (\ref{eq:#2})}
\newcommand{\Erefss}[2]{Equations~(\ref{eq:#1}) - (\ref{eq:#2})}
\newcommand{\trefss}[2]{Tables~\ref{tab:#1} - \ref{tab:#2}}
%misc.
%Note: tarticle has ``several'' changes from article
%in this vein.
% some simplifying commands
\newcommand{\eg}{{\it e.g.}}
\newcommand{\ie}{{\it i.e.}}
\newcommand{\hisparcbox}[1]{%
    \path[hisparc,#1]
        (-.4, .7) .. controls (0, .75) ..  (.4, .7) --
        (.35, -1.7) ..  controls(0, -1.72) ..  (-.35, -1.7) -- cycle;
}
% derivative - d
\newcommand{\ud}{\,\mathrm{d}}

%Get rid of navigation icons
\setbeamertemplate{navigation symbols}{}
\useoutertheme{infolines}


%Aggie-themed
\pgfdeclareimage[height=0.1in]{TAMUlogo}{tamu_engineering.png}
\logo{\raisebox{-8pt}{\pgfuseimage{TAMUlogo}}}
\titlegraphic{\includegraphics[height=0.18\textheight]{tamu_seal.png}}
%Michigan-themed
%\pgfdeclareimage[height=0.1in]{UMlogo}{michigan_engineering.png}
%\logo{\raisebox{-8pt}{\pgfuseimage{UMlogo}}}
%\titlegraphic{\includegraphics[height=0.2\textheight]{michigan_block_m.png}}


%%%%%%%%%%%%%%%%%%%%%%%%%%%%%%%%%%%%%%%%%%%%%%%%%%%%%%%%%%%%%%%
% Optional packages, used to show off certain tricks

\newlength \figwidth
\setlength \figwidth {0.5\textwidth}

\setlength{\leftmargin}{-2cm}
\setlength{\rightmargin}{-2cm}

%%%%%%%%%%%%%%%%%%%%%%%%%%%%%%%%%%%%%%%%%%%%%%%%%%%%%%%%%%%%%%%

\mode<presentation>
{
\usepackage[english]{babel}
\usetheme{Frankfurt}

%Make it Aggie Maroon
\usecolortheme[RGB={80,0,0}]{structure}  
%Or Michigan Blue
%\usecolortheme[RGB={0,0,153}]{structure}  
%Or Michigan Maize
%\usecolortheme[RGB={255,204,0}]{structure}  

  % This will typeset only the frames (or slides) that have the given label ("current" in this case).
%  \includeonlyframes{current}
}

\title{Final Project Results}
\subtitle{Acceleration of non-LTE calculations via parallel cyclic reduction}

\author[Holladay]{\large Daniel~A.~Holladay}

%TAMU
\institute[TAMU]{Department of Nuclear Engineering\\
Texas A\&M University \\
College Station, TX 77843-3133\\[1ex]
\href{mailto:dholladay00@tamu.edu}{dholladay00@tamu.edu}}

%U of M
%\institute[Michigan]{Department of Nuclear Engineering and Radiological Sciences\\
%University of Michigan \\
%Ann Arbor, MI 48109 \\[1ex]
%\href{mailto:andthoti@umich.edu}{andthoti@umich.edu}}

\date{\today}

% You can override the default acknowledgment, and address if you want
%\acknowledgement{*Submitted in partial fulfillment of the requirements of NUEN 610 \\
%(Nuclear Reactor Design)}
%\address{Nuclear Engineering Department \\
%            Texas A\&M University \\
%            College Station, TX 77843-3133}}

% If you don't want the menu section outline above the title, do this:
%\setbeamertemplate{headline}{}

%%%%%%%%%%%%%%%%%%%%%%%%%%%%%%%%%%%%%%%%%%%%%%%%%%%%%%%%%%%%%%%%%%%%%%%%%%%%%%%%%%%%%%%%%%%%%
\begin{document}

%%%%%%%%%%%%%%%%%%%%%%%%%%%%%%%%%%%%%%%%%%%%%%%%%%%%%%%%%%%%%%%%%%%%%%%%%%%%%%%%%%%%%%%%%%%%%
%  All this typeout stuff simply gets printed to the screen as the document
% is compiled.  It helps get stuff working
\typeout{***********************************************************************************}
\typeout{titlepage}

\begin{frame}[label=title,plain]
\titlepage
\end{frame}

\begin{frame}
\frametitle{Parallel algorithm}

\begin{center}
\begin{tikzpicture}
    [ x=.5mm, y=.5mm,
      font={\sffamily},
      station/.style={fill=gray},
      hut/.style={fill=lightgray},
      cluster/.style={fill=yellow!30, rounded corners=2pt},
      road/.style={fill=blue!10},
      calorimeter/.style={fill=green!30},
      tracker/.style={fill=red!30},
      hisparc/.style={fill=red, rounded corners=.15pt},
      hisparcgps/.style={fill=red},
      axis/.style={gray,very thick,->,>=stealth'},
      ruler/.style={gray,|<->|,>=stealth'},
      spy using outlines={red, circle, magnification=4, size=25 * 4,
                          connect spies}
    ]
    % Clusters
%    \foreach \i in {-1.5, -0.5, ..., 1.5} {
%        \foreach \j in {-1.5, -0.5, ..., 1.5} {
%            \path[cluster, shift={(\i * 52, \j * 52)}]
%                (-22, -22) rectangle (22, 22);
%        }
%    }

    % Roads
%    \foreach \i in {-1.5, -0.5, ..., 1.5} {
%        \path[road, shift={(0, \i * 52 + 2)}]
%            (-105, -1.5) rectangle (105, 1.5);
%    }

    % Detector array
    \foreach \i in {-7.5, -6.5, ..., 7.5} {
        \foreach \j in {-7.5, -6.5, ..., 7.5} {
            \path[station, shift={(\i * 10, \j * 10)}]
                (-1.2, -1.2) rectangle (1.2, 1.2);
        }
    }
    % Two detectors are slightly displaced
%    \path[station, shift={(-6.5, 23.5)}] (-1.2, -1.2) rectangle (1.2, 1.2);
%    \path[station, shift={(6.5, 23.5)}] (-1.2, -1.2) rectangle (1.2, 1.2);
%
%    % Electronic huts
%    \foreach \i in {-1.5, -0.5, ..., 1.5} {
%        \foreach \j in {-1.5, -0.5, ..., 1.5} {
%            \path[hut, shift={(\i * 52, \j * 52 - 1.5)}]
%                (-2.5, -1.2) rectangle (2.5, 1.2);
%        }
%    }
%
%    % Central detector
%    \fill[white] (-13, -13) rectangle (13, 21);
%    \path[calorimeter, shift={(0, 5)}] (-10, -8) rectangle (10, 8);
%
%    % Muon tracker
%    \path[tracker, shift={(0, 4)}] (-2.7, 32.5) rectangle (2.7, 70.5);
%
%    % HiSPARC station
%    \hisparcbox{shift={(65.0, 15.05)}, rotate=180}
%    \hisparcbox{shift={(65.0, 20.82)}, rotate=180}
%    \hisparcbox{shift={(70.0, 23.71)}, rotate=-90}
%    \hisparcbox{shift={(60.0, 23.71)}, rotate=90}
%    \path[hisparcgps, shift={(65.0, 23.71)}] (0, 0) circle (.3);
%
%    % Draw rulers
%    \draw[ruler] (-97.5, -105) -- (97.5, -105)
%        node[midway,below] {\SI{195}{\meter}};
%    \draw[ruler] (84.5, 105) -- (97.5, 105)
%        node[midway,above] {\SI{13}{\meter}};
%    \draw[ruler] (-105, -97.5) -- (-105, 97.5)
%        node[midway,above,sloped] {\SI{195}{\meter}};
%    \draw[ruler] (105, 84.5) -- (105, 97.5)
%        node[midway,below,sloped] {\SI{13}{\meter}};
%
%    % Spy
%    \spy on (65.0, 20.82) in node at (30, -50);

\end{tikzpicture}
\end{center}

\end{frame}

\end{document}
